% =====================================================================
%  Stoff und Vlies im Stickrahmen, Schnitt durch eine Rahmenseite.
%  Schematisch, nicht maßstäblich. Selbst gezeichnet.
% =====================================================================
\begin{tikzpicture}[x=2.4mm, y=2.4mm, font=\scriptsize]
  % Außenrahmen und Innenrahmen im Schnitt
  \draw[rahmen] (0,0) rectangle (5,9);
  \draw[rahmen, fill=black!35] (6,-0.9) rectangle (11,8.4);
  % Vlies (unten) und Stoff (oben) als Lagen
  \draw[line width=1.6pt, black!35, line join=round]
    (-6,6.3) -- (-0.6,9.4) -- (5.3,9.4) -- (5.3,-1.3) -- (40,-1.3);
  \draw[line width=1.6pt, stoffrand, line join=round]
    (-6,7.1) -- (-0.6,9.9) -- (5.75,9.9) -- (5.75,-0.9) -- (40,-0.9);
  \draw[hilfslinie, dashed] (-8,0) -- (0,0);
  \draw[hilfslinie, dashed] (11,-0.9) -- (16,-0.9);

  % Nadel über der Stofffläche
  \draw[metall, fill=black!55] (31.6,12) -- (32.4,12) -- (32.2,2) -- (32,1) -- (31.8,2) -- cycle;

  \txt[anchor=south]{2.5,15}{2.5,9}{Außenrahmen\\(Schraube)}
  \txt[anchor=south]{12,15}{9,8.4}{Innenrahmen}
  \txt[anchor=west]{18,7}{20,-0.6}{Stoff, rechte Seite oben}
  \txt[anchor=west]{18,-6}{22,-1.5}{Vlies darunter}
  \txt[anchor=east]{-8,-4}{-4.5,6.9}{Stoff und Vlies\\stehen außen über}
  \txt[anchor=north]{8.5,-6}{8.5,-0.9}{Unterseite des Innenrahmens nicht\\höher als die des Außenrahmens}
  \node[anchor=south] at (32,12.5) {Nadel};
\end{tikzpicture}
